
\phantomsection
\section*{Глава 2. Проектирование веб-приложения}
\addcontentsline{toc}{section}{Глава 2. Проектирование веб-приложения}
\stepcounter{section}

\subsection{Анализ целевой аудитории}

Разрабатываемое веб-приложение ориентировано на пользователей, которые участвуют в процессе оказания медицинских услуг в частной клинике. Для корректного проектирования интерфейса и разграничения доступа выделены три основные роли: пациент, врач и администратор. Дополнительно система должна учитывать неавторизованного посетителя, который может ознакомиться с общими сведениями о клинике, но не имеет доступа к медицинским данным.

{\small
\setlength{\tabcolsep}{4pt}
\renewcommand{\arraystretch}{1.12}
\begin{longtable}{|>{\raggedright\arraybackslash}p{0.24\textwidth}|>{\raggedright\arraybackslash}p{0.30\textwidth}|>{\raggedright\arraybackslash}p{0.36\textwidth}|}
\caption{Характеристика целевой аудитории веб-приложения}\label{tab:audience}\\
\hline
\textbf{Роль} & \textbf{Основные потребности} & \textbf{Типовые действия в системе} \\ \hline
\endfirsthead
\tabcont{\ref{tab:audience}}
\textbf{Роль} & \textbf{Основные потребности} & \textbf{Типовые действия в системе} \\ \hline
\endhead
Неавторизованный посетитель & Получение общей информации о клинике, услугах и врачах. & Просмотр публичных страниц, переход к регистрации или форме входа. \\ \hline
Пациент & Быстрая запись на приём, просмотр расписания, доступ к истории посещений и медицинским документам. & Регистрация, вход в личный кабинет, выбор врача и услуги, создание или отмена записи, просмотр рекомендаций врача. \\ \hline
Врач & Удобный просмотр расписания, доступ к медицинской карте пациента, фиксация результатов приёма. & Просмотр списка приёмов, открытие карты пациента, добавление жалоб, диагноза, назначений и файлов. \\ \hline
Администратор & Управление расписанием клиники, контроль записей, поддержка справочников и пользователей. & Создание врачей и услуг, настройка графика работы, подтверждение или перенос записей, управление учётными записями. \\ \hline
\end{longtable}
}

Для пациента интерфейс должен быть простым и понятным: основные действия должны выполняться за минимальное количество шагов. Для врача важны быстрый доступ к расписанию и медицинской карте пациента. Для администратора требуется более полный набор инструментов, так как эта роль отвечает за поддержание актуальности данных в системе.

\subsection{Описание функциональности приложения}

Функциональность приложения определяется через действия пользователей и системные ограничения. Общая диаграмма вариантов использования представлена на рисунке~\ref{fig:use_case}. Диаграмма отражает связь между тремя основными ролями и ключевыми действиями: авторизацией, записью на приём, управлением расписанием, ведением медицинских карт и просмотром истории посещений.

\begin{figure}[ht!]
\centering
\includegraphics[height=1\textwidth]{use_case_diagram}
\caption{Диаграмма вариантов использования веб-приложения}
\label{fig:use_case}
\end{figure}

Ключевые пользовательские сценарии можно описать с помощью User Story:

\begin{enumerate}
\item Я, как пациент, хочу выбрать врача, услугу и свободное время, чтобы самостоятельно записаться на приём без звонка в клинику.
\item Я, как пациент, хочу видеть историю приёмов и рекомендации врача, чтобы иметь доступ к медицинской информации после посещения.
\item Я, как врач, хочу просматривать своё расписание на день, чтобы заранее видеть список пациентов и время приёма.
\item Я, как врач, хочу заполнять медицинскую карту пациента, чтобы фиксировать результаты приёма и назначения.
\item Я, как администратор, хочу управлять расписанием врачей, услугами и записями, чтобы обеспечивать корректную работу клиники.
\item Я, как администратор, хочу иметь возможность выгружать или прикреплять документы, чтобы поддерживать актуальность медицинской информации пациента.
\end{enumerate}

На основании анализа сценариев сформирован перечень функциональных требований, представленный в таблице~\ref{tab:functional_requirements}.

{\small
\setlength{\tabcolsep}{4pt}
\renewcommand{\arraystretch}{1.12}
\begin{longtable}{|>{\raggedright\arraybackslash}p{0.13\textwidth}|>{\raggedright\arraybackslash}p{0.52\textwidth}|>{\raggedright\arraybackslash}p{0.25\textwidth}|}
\caption{Функциональные требования к веб-приложению}\label{tab:functional_requirements}\\
\hline
\textbf{Номер} & \textbf{Требование} & \textbf{Роль} \\ \hline
\endfirsthead
\tabcont{\ref{tab:functional_requirements}}
\textbf{Номер} & \textbf{Требование} & \textbf{Роль} \\ \hline
\endhead
1 & Система должна обеспечивать регистрацию, вход и выход пользователей. & Пациент, врач, администратор \\ \hline
2 & Система должна разграничивать доступ к страницам и операциям по ролям. & Пациент, врач, администратор \\ \hline
3 & Пациент должен иметь возможность просматривать список врачей, услуг и свободных временных интервалов. & Пациент \\ \hline
4 & Пациент должен иметь возможность создать, просмотреть или отменить запись на приём. & Пациент \\ \hline
5 & Врач должен иметь возможность просматривать собственное расписание и карточку приёма. & Врач \\ \hline
6 & Врач должен иметь возможность заполнять медицинскую карту пациента по итогам приёма. & Врач \\ \hline
7 & Администратор должен управлять врачами, услугами, расписанием, пользователями и записями. & Администратор \\ \hline
8 & Система должна поддерживать загрузку, просмотр и выгрузку медицинских документов. & Врач, пациент, администратор \\ \hline
9 & Система должна отображать пациенту личный кабинет с будущими записями, историей посещений и медицинской картой. & Пациент \\ \hline
\end{longtable}
}

Кроме функциональных требований, для приложения важны нефункциональные требования: защита персональных данных, устойчивость к ошибочным действиям пользователя, адаптивность интерфейса, понятная навигация и возможность резервного копирования базы данных.

\subsection{Проектирование модели данных}

Модель данных должна отражать связи между основными объектами предметной области. Ключевыми сущностями являются пользователь, врач, пациент, услуга, расписание, запись на приём, медицинская карта, посещение и медицинский документ. Упрощённая ER-диаграмма представлена на рисунке~\ref{fig:er_diagram}.

\begin{figure}[ht!]
\centering
\includegraphics[width=0.90\textwidth]{er_diagram}
\caption{Упрощённая ER-диаграмма базы данных}
\label{fig:er_diagram}
\end{figure}

Логическая схема базы данных строится на следующих правилах: один пользователь имеет одну роль; пациент связан с одной учётной записью; врач связан с одной учётной записью и может иметь множество слотов расписания; запись на приём связывает пациента, врача, услугу и свободный временной интервал; медицинская карта принадлежит пациенту и содержит историю посещений; документы могут быть связаны с пациентом и конкретным посещением.

{\small
\setlength{\tabcolsep}{4pt}
\renewcommand{\arraystretch}{1.12}
\begin{longtable}{|>{\raggedright\arraybackslash}p{0.28\textwidth}|>{\raggedright\arraybackslash}p{0.22\textwidth}|>{\raggedright\arraybackslash}p{0.40\textwidth}|}
\caption{Логическая схема данных}\label{tab:logical_schema}\\
\hline
\textbf{Связь} & \textbf{Кратность} & \textbf{Описание} \\ \hline
\endfirsthead
\tabcont{\ref{tab:logical_schema}}
\textbf{Связь} & \textbf{Кратность} & \textbf{Описание} \\ \hline
\endhead
users - patients & 1:1 & Учётная запись пациента содержит данные для входа, а профиль пациента хранит персональную информацию. \\ \hline
users - doctors & 1:1 & Учётная запись врача используется для входа и разграничения доступа к расписанию и медицинским картам. \\ \hline
doctors - schedule\_slots & 1:N & Один врач имеет множество временных интервалов расписания. \\ \hline
patients - appointments & 1:N & Один пациент может иметь несколько записей на приём. \\ \hline
appointments - visits & 1:0..1 & После завершения приёма запись может быть связана с результатом посещения. \\ \hline
patients - medical\_records & 1:1 & Медицинская карта создаётся для конкретного пациента. \\ \hline
visits - documents & 1:N & К одному посещению могут быть прикреплены результаты анализов, заключения или иные файлы. \\ \hline
\end{longtable}
}

Основные таблицы базы данных представлены в таблице~\ref{tab:database_tables}. В рамках курсового проекта такая структура достаточна для реализации базовой логики приложения и может быть расширена при дальнейшем развитии системы.

{\small
\setlength{\tabcolsep}{4pt}
\renewcommand{\arraystretch}{1.12}
\begin{longtable}{|>{\raggedright\arraybackslash}p{0.20\textwidth}|>{\raggedright\arraybackslash}p{0.34\textwidth}|>{\raggedright\arraybackslash}p{0.36\textwidth}|}
\caption{Основные таблицы базы данных}\label{tab:database_tables}\\
\hline
\textbf{Таблица} & \textbf{Ключевые поля} & \textbf{Назначение} \\ \hline
\endfirsthead
\tabcont{\ref{tab:database_tables}}
\textbf{Таблица} & \textbf{Ключевые поля} & \textbf{Назначение} \\ \hline
\endhead
users & id, email, password\_hash, role, created\_at & Хранение учётных записей и роли пользователя. \\ \hline
patients & id, user\_id, full\_name, birth\_date, phone & Данные пациента, связанные с учётной записью. \\ \hline
doctors & id, user\_id, full\_name, specialization, cabinet & Данные врача, специальность и кабинет приёма. \\ \hline
services & id, title, description, price, duration\_minutes & Справочник медицинских услуг. \\ \hline
schedule\_slots & id, doctor\_id, date, start\_time, end\_time, status & Доступные временные интервалы врача. \\ \hline
appointments & id, patient\_id, doctor\_id, service\_id, slot\_id, status & Записи пациентов на приём и их статусы. \\ \hline
medical\_records & id, patient\_id, created\_at, updated\_at & Основная медицинская карта пациента. \\ \hline
visits & id, record\_id, appointment\_id, complaints, diagnosis, recommendations & Результаты конкретного посещения врача. \\ \hline
documents & id, patient\_id, visit\_id, file\_path, file\_type & Прикреплённые медицинские файлы и результаты исследований. \\ \hline
\end{longtable}
}

Физическая структура одной из ключевых таблиц - таблицы записей на приём - представлена в таблице~\ref{tab:appointments}. Она показывает, какие типы данных и ограничения должны быть использованы при реализации базы данных в PostgreSQL.

{\small
\setlength{\tabcolsep}{4pt}
\renewcommand{\arraystretch}{1.12}
\begin{longtable}{|>{\raggedright\arraybackslash}p{0.20\textwidth}|>{\raggedright\arraybackslash}p{0.18\textwidth}|>{\raggedright\arraybackslash}p{0.26\textwidth}|>{\raggedright\arraybackslash}p{0.26\textwidth}|}
\caption{Структура таблицы записей на приём (appointments)}\label{tab:appointments}\\
\hline
\textbf{Поле} & \textbf{Тип} & \textbf{Ограничение} & \textbf{Описание} \\ \hline
\endfirsthead
\tabcont{\ref{tab:appointments}}
\textbf{Поле} & \textbf{Тип} & \textbf{Ограничение} & \textbf{Описание} \\ \hline
\endhead
id & SERIAL & PRIMARY KEY & Уникальный идентификатор записи. \\ \hline
patient\_id & INTEGER & FOREIGN KEY & Ссылка на пациента. \\ \hline
doctor\_id & INTEGER & FOREIGN KEY & Ссылка на врача. \\ \hline
service\_id & INTEGER & FOREIGN KEY & Ссылка на выбранную услугу. \\ \hline
slot\_id & INTEGER & FOREIGN KEY, UNIQUE & Ссылка на временной интервал расписания; ограничение UNIQUE предотвращает двойную запись на один слот. \\ \hline
status & VARCHAR(20) & NOT NULL & Статус записи: создана, подтверждена, отменена, завершена. \\ \hline
created\_at & TIMESTAMP & NOT NULL & Дата и время создания записи. \\ \hline
\end{longtable}
}

Также необходимо предусмотреть хранение файлов. В базе данных должен сохраняться не сам файл, а путь к нему, тип файла, владелец документа и связь с посещением. Такой подход упрощает работу с базой данных и позволяет организовать отдельное хранилище загруженных медицинских документов.

\subsection{Разработка макетов страниц}

На этапе проектирования интерфейса определены основные страницы приложения:

\begin{enumerate}
\item Публичная главная страница с описанием клиники, услуг и переходом к записи;
\item Страницы регистрации и входа;
\item Личный кабинет пациента с будущими записями, историей посещений и медицинской картой;
\item Страница выбора врача, услуги, даты и времени приёма;
\item Кабинет врача со списком приёмов и формой заполнения медицинской карты;
\item Административная панель для управления врачами, услугами, расписанием и записями.
\end{enumerate}

Пример структуры личного кабинета пациента представлен на рисунке~\ref{fig:wireframe_patient}.

\begin{figure}[ht!]
\centering
\includegraphics[width=0.80\textwidth]{wireframe_patient}
\caption{Макет структуры личного кабинета пациента}
\label{fig:wireframe_patient}
\end{figure}

При проектировании интерфейса используются следующие принципы:

\begin{enumerate}
\item Основная навигация должна быть доступна с любой страницы после входа в систему;
\item Формы должны содержать минимально необходимый набор полей;
\item Таблицы расписания и записей должны поддерживать фильтрацию по дате, врачу и статусу;
\item Действия, связанные с изменением медицинской карты, должны быть доступны только врачу и администратору;
\item Пациент не должен видеть данные других пациентов;
\item После успешного действия пользователь должен получать понятное уведомление.
\end{enumerate}

Таким образом, проектирование веб-приложения включает выделение целевой аудитории, описание функциональности через диаграмму вариантов использования и пользовательские истории, разработку модели данных, логической и физической схемы БД, а также подготовку макетов страниц. Эти результаты являются основой для последующей разработки программной реализации.
